\documentclass[11pt,a4paper]{amsart}

\usepackage[utf8]{inputenc}
\usepackage[T1]{fontenc}
\usepackage{lmodern}
\usepackage[english]{babel}
\usepackage{amsmath, amssymb, amsthm}
\usepackage{hyperref}
\usepackage{enumerate}
\usepackage{geometry}
\usepackage{float}
\usepackage{booktabs}
\usepackage{csquotes}

\theoremstyle{plain}
\newtheorem{theorem}{Theorem}[section]
\newtheorem{lemma}[theorem]{Lemma}
\newtheorem{corollary}[theorem]{Corollary}
\newtheorem{conjecture}[theorem]{Conjecture}

\theoremstyle{definition}
\newtheorem{definition}[theorem]{Definition}

\theoremstyle{remark}
\newtheorem{remark}[theorem]{Remark}

\usepackage[
    backend=biber,
    style=numeric,
    sorting=none
]{biblatex}
\addbibresource{bibliography.bib}

\geometry{margin=2.8cm}

\title{A Linear Algorithm for Huang's Quadratic Form\\
on Numerical Semigroups and Density Asymptotics}

\author{Artur Flamandzki}
\thanks{E-mail: \texttt{flamandzki.artur@gmail.com}.}

\date{\today}

\begin{document}

\begin{abstract}
Let $G = \mathbb{N} \setminus \langle a, b \rangle$ be the gap set
of the numerical semigroup generated by coprime $a < b$, and $N = |G|$.
Huang~\cite{Huang2026} defined the quadratic form
$Q(\mathbf{n}) = \sum_{i,j} K(j-i)\,n_i n_j$ on $\mathbb{R}^G$, where
$K(d) = \mathbf{1}_{d \geq 0} - \mathbf{1}_{d \geq a} - \mathbf{1}_{d \geq b}
+ \mathbf{1}_{d \geq a+b}$, and showed that it recovers the
$\mathrm{dinv}$ statistic on rational Dyck paths.
The naive evaluation of $Q$ requires $O(N^2)$ operations.

We prove that the interval structure of $K$ allows one to reduce $Q$
to a linear combination of sliding-window sums, yielding an $O(N)$
algorithm (Theorem~\ref{thm:linear}). We generalise to semigroups
$\langle p_1,\ldots,p_n \rangle$ with $n$ generators, where the
generalised kernel $K^{(n)}$ defined by inclusion--exclusion
--- with $2^n$ terms --- has at most $w(n) \leq \sigma_n$
active windows due to parity cancellation;
these are computable incrementally in time $O(n \cdot \sigma_n)$
(Theorem~\ref{thm:general}).
Empirically $w(n) \sim \sigma_n$ while $w(n) \ll 2^n$:
exponential combinatorics reduces to a sparse interval structure.
For a sequence of pairwise distinct integers $p_i \geq 2$ satisfying
$\log p_n = o(n)$ we prove
\[
  \frac{w(n)}{2^n - 1} \;=\; o\!\left(\,\prod_{i=1}^{n}
    \Bigl(1 - \tfrac{1}{p_i}\Bigr)\right), \qquad n \to \infty
\]
(Theorem~\ref{thm:density}).
\end{abstract}

\maketitle

% ==================================================
\section{Introduction}
\label{sec:intro}
% ==================================================

\subsection*{1.1. Background and motivation}

Let $a, b \in \mathbb{N}$, $\gcd(a,b) = 1$, $a < b$.
The numerical semigroup generated by $a$ and $b$ is
$\langle a, b \rangle = \{ xa + yb : x, y \in \mathbb{Z}_{\geq 0} \}$.
By the Sylvester--Frobenius theorem the gap set
$G = \mathbb{N} \setminus \langle a, b \rangle$ is finite,
$|G| = N = \tfrac{(a-1)(b-1)}{2}$.
We write $G = \{g_0 < g_1 < \cdots < g_{N-1}\}$.

Huang~\cite{Huang2026} defined on $\mathbb{R}^G$ the quadratic form
\[
  Q(\mathbf{n}) = \sum_{i,j \in G} K(j - i)\, n_i n_j,
  \qquad
  K(d) = \mathbf{1}_{d \geq 0} - \mathbf{1}_{d \geq a}
         - \mathbf{1}_{d \geq b} + \mathbf{1}_{d \geq a+b},
\]
and showed that when $\mathbf{n} = \mathbf{1}_D$ is the indicator vector
of a Young subdiagram $D \subseteq G$, then $Q(\mathbf{1}_D) = \mathrm{dinv}(D)$,
where $\mathrm{dinv}$ is the Gorsky--Mazin statistic on rational Dyck paths.
An independent formal verification of this theorem in
Lean~4/Mathlib is given in~\cite{AxiomProver2026}.

The paper~\cite{Huang2026} and its formalisation~\cite{AxiomProver2026}
concern the identity $Q(\mathbf{1}_D) = \mathrm{dinv}(D)$.
The present paper concerns the computational efficiency of $Q$
and its generalisation to $n$ generators.

Direct evaluation of $Q$ from its definition requires $N^2$ operations.
The key observation is that $K$ takes at most three nonzero values
on five disjoint intervals (Lemma~\ref{lem:K-intervals}), which reduces
the computation of $Q$ to two sliding-window sums, giving an $O(N)$
algorithm.

The generalisation to $n$ generators defines the kernel $K^{(n)}$
by inclusion--exclusion over all $2^n$ subsets of $\{p_1, \ldots, p_n\}$.
The interval structure of $K^{(n)}$ is preserved, but the number of
active windows $w(n)$ --- intervals on which $K^{(n)} \neq 0$ ---
is far smaller than $2^n$ due to parity cancellation.
We prove $w(n) \leq \sigma_n$ and give an incremental algorithm
computing $w(n)$ and the active windows themselves.

The ratio $w(n)/(2^n - 1)$ is naturally linked to the product
$P(n) = \prod_{i=1}^n (1 - 1/p_i)$: under the condition
$\log p_n = o(n)$ we show $w(n)/(2^n-1) = o(P(n))$,
i.e., the density of active windows is asymptotically negligible
relative to $P(n)$.

The generalisation to $n$ generators does not address the Frobenius
problem, which asks for an explicit formula for the largest element of
$\mathbb{Z}_{\geq 0} \setminus \langle p_1,\ldots,p_n \rangle$
and remains open for $n \geq 3$~\cite{Ramirez2005}.
The kernel $K^{(n)}$ is defined by inclusion--exclusion over the
generator set; the algorithm takes $G$ as input.

\subsection*{1.2. Definitions and main results}

\textbf{Two generators.}
For fixed $k$ and $\mathbf{n} = (n_j)_{j \in G} \in \mathbb{R}^G$
define the window sums
\begin{align*}
  W^+_k &:= \sum_{\substack{j \in G \\ g_k \leq g_j < g_k + a}} n_j, \qquad
  W^-_k := \sum_{\substack{j \in G \\ g_k + b \leq g_j < g_k + a + b}} n_j.
\end{align*}

\begin{theorem}[Linear algorithm, two generators]
\label{thm:linear}
For every $\mathbf{n} \in \mathbb{R}^G$,
\[
  Q(\mathbf{n}) = \sum_{k=0}^{N-1} n_k \bigl(W^+_k - W^-_k\bigr).
\]
Both families $(W^+_k)$ and $(W^-_k)$ are jointly computable
in time $O(N)$ and space $O(N)$ using four sliding-window
pointers on $G$.
\end{theorem}

Theorem~\ref{thm:linear} reduces the cost from $O(N^2)$ to $O(N)$.
Since $N = \tfrac{(a-1)(b-1)}{2}$, the algorithm runs in time
$O(ab)$ versus $O(a^2 b^2)$ for the naive approach.

\textbf{Multiple generators.}
Let $\mathbf{p} = (p_1, \ldots, p_n)$, $p_i \geq 2$ pairwise distinct,
$\gcd(p_1,\ldots,p_n) = 1$, $\sigma_n = \sum_{i=1}^n p_i$.
The generalised kernel
\[
  K^{(n)}(d) := \sum_{S \subseteq \{p_1,\ldots,p_n\}} (-1)^{|S|}
               \mathbf{1}_{d \geq \sigma(S)}, \qquad
  \sigma(S) = \sum_{p \in S} p, \quad \sigma(\emptyset) = 0,
\]
is nonzero on at most $w(n)$ intervals between consecutive
breakpoints; we write $w(n)$ for the number of active windows.

\begin{theorem}[Sliding-window algorithm, $n$ generators]
\label{thm:general}
Let $G = \mathbb{N} \setminus \langle p_1,\ldots,p_n \rangle$.
The form $Q(\mathbf{n}) = \sum_{i,j \in G} K^{(n)}(j-i)\,n_i n_j$
is computable in time $O(T_{\mathrm{prep}} + w(n) \cdot N)$, where
$T_{\mathrm{prep}} = O(n \cdot \sigma_n)$.
One has $w(n) \leq \sigma_n$.
The set of active windows and the values of $K^{(n)}$ on them are
computable incrementally: the transition from $n-1$ to $n$ generators
costs $O(\sigma_{n-1})$.
\end{theorem}

\textbf{Density asymptotics.}
For a sequence of pairwise distinct integers $p_i \geq 2$
define the density product and density defect:
\[
  P(n) := \prod_{i=1}^{n} \!\left(1 - \frac{1}{p_i}\right), \qquad
  \delta(n) := P(n) - \frac{w(n)}{2^n - 1}.
\]

\begin{remark}
\label{rem:terminology}
The term \emph{gap} refers exclusively to an element of
$G = \mathbb{N} \setminus \langle p_1,\ldots,p_n \rangle$.
The symbol $\delta(n)$ refers exclusively to the quantity $P(n) - w(n)/(2^n-1)$.
The two notions are independent.
\end{remark}

\begin{remark}
\label{rem:mertens}
The product $P(n) = \prod_{i=1}^n (1 - 1/p_i)$ is defined
for any sequence of pairwise distinct integers $p_i \geq 2$
and does not require the $p_i$ to be prime.
When the $p_i$ are consecutive primes, $P(n)$ coincides with
the classical Mertens product; in that case $P(n) \to 0$.
\end{remark}

\begin{theorem}[Density asymptotics]
\label{thm:density}
Let $(p_i)_{i \geq 1}$ be a sequence of pairwise distinct
integers $\geq 2$ satisfying $\log p_n = o(n)$.
Then
\[
  \frac{w(n)}{2^n - 1} = o(P(n)), \qquad n \to \infty,
\]
equivalently, $\delta(n)/P(n) \to 1$.
\end{theorem}

The condition $\log p_n = o(n)$ is satisfied by any family
of prime generators ($p_n \sim n \log n$), twin primes,
arithmetic progressions of primes, and sequences of random
odd numbers of prime density.

% ==================================================
\section{Preliminaries}
\label{sec:prelim}
% ==================================================

\subsection*{Two-generator kernel}

\begin{lemma}[Interval structure of $K$]
\label{lem:K-intervals}
For all $d \in \mathbb{Z}$:
\[
  K(d) =
  \begin{cases}
     0  & d < 0,            \\
    +1  & 0 \leq d < a,     \\
     0  & a \leq d < b,     \\
    -1  & b \leq d < a+b,   \\
     0  & d \geq a+b.
  \end{cases}
\]
\end{lemma}

\begin{proof}
Direct inspection of the four indicators $\mathbf{1}_{d \geq 0}$,
$-\mathbf{1}_{d \geq a}$, $-\mathbf{1}_{d \geq b}$,
$+\mathbf{1}_{d \geq a+b}$ for $a < b$.
The five cases are mutually exclusive and exhaustive.
\end{proof}

\begin{corollary}
\label{cor:support}
$K(d) \neq 0$ implies $0 \leq d < a+b$.
The sum $\sum_{d \in \mathbb{Z}} K(d) = a - (b - a) = 2a - b$.
\end{corollary}

\subsection*{Multi-generator kernel}

\begin{definition}[Parity function]
\label{def:parity}
For $\mathbf{p} = (p_1,\ldots,p_n)$ and $s \in \mathbb{Z}_{\geq 0}$:
\[
  \pi_n(s) := \sum_{\substack{S \subseteq \{p_1,\ldots,p_n\} \\
                              \sigma(S) = s}} (-1)^{|S|}.
\]
One has $K^{(n)}(d) = \sum_{s \leq d} \pi_n(s)$,
so $K^{(n)}$ is constant on intervals between consecutive
$s$ with $\pi_n(s) \neq 0$.
\end{definition}

\begin{lemma}[Vanishing outside $[0, \sigma_n)$]
\label{lem:K-vanish}
$K^{(n)}(d) = 0$ for $d < 0$ and for $d \geq \sigma_n$.
\end{lemma}

\begin{proof}
For $d < 0$: every indicator $\mathbf{1}_{d \geq \sigma(S)} = 0$.
For $d \geq \sigma_n$: all indicators equal $1$, and
$\sum_{S \subseteq \{p_1,\ldots,p_n\}} (-1)^{|S|} = (1-1)^n = 0$
by the binomial theorem.
\end{proof}

\begin{lemma}[Parity recurrence]
\label{lem:parity-recurrence}
\[
  \pi_{n+1}(s) = \pi_n(s) - \pi_n(s - p_{n+1}).
\]
\end{lemma}

\begin{proof}
Subsets of $\{p_1,\ldots,p_{n+1}\}$ with sum $s$ split into:
subsets $S \subseteq \{p_1,\ldots,p_n\}$ with $\sigma(S) = s$,
contributing $\pi_n(s)$, and subsets $S \cup \{p_{n+1}\}$
with $\sigma(S) = s - p_{n+1}$, contributing $-\pi_n(s - p_{n+1})$.
\end{proof}

% ==================================================
\section{Proof of Theorem~\ref{thm:linear}}
\label{sec:proof-linear}
% ==================================================

\subsection*{3.1. Reduction of the quadratic form}

\begin{lemma}
\label{lem:inner-sum}
For fixed $i \in G$:
\[
  \sum_{j \in G} K(j - i)\, n_j = W^+_i - W^-_i.
\]
\end{lemma}

\begin{proof}
By Lemma~\ref{lem:K-intervals}: $K(j - i) = +1$ when
$g_i \leq g_j < g_i + a$, $K(j-i) = -1$ when
$g_i + b \leq g_j < g_i + a + b$, and $K(j-i) = 0$ otherwise.
The definitions of $W^+_i$ and $W^-_i$ correspond precisely
to these two intervals.
\end{proof}

\begin{proof}[Proof of Theorem~\ref{thm:linear}]
Rearranging the sum:
\[
  Q(\mathbf{n}) = \sum_{i \in G} n_i \sum_{j \in G} K(j-i)\,n_j
  = \sum_{i \in G} n_i (W^+_i - W^-_i),
\]
where the first equality follows by symmetry (index swap)
and the second from Lemma~\ref{lem:inner-sum}.

Linearity: maintain four pointers
$\mathrm{lo}^+, \mathrm{hi}^+, \mathrm{lo}^-, \mathrm{hi}^-$
into sorted $G$, initialised to $0$.
For $k = 0, 1, \ldots, N-1$:
advance $\mathrm{hi}^+$ while $g_{\mathrm{hi}^+} < g_k + a$;
advance $\mathrm{lo}^+$ while $g_{\mathrm{lo}^+} < g_k$;
advance $\mathrm{hi}^-$ while $g_{\mathrm{hi}^-} < g_k + a + b$;
advance $\mathrm{lo}^-$ while $g_{\mathrm{lo}^-} < g_k + b$.
Each pointer is non-decreasing and traverses $G$ at most once:
at most $4N$ updates in total.
\end{proof}

\begin{corollary}
\label{cor:speedup}
For $(a, b) = (97, 101)$, $N = 4801$: the $O(N)$ algorithm performs
$19\,204$ operations versus $23\,049\,601$ for $O(N^2)$;
ratio $\approx 1200$.
In general $N^2 / (4N) = N/4 = (a-1)(b-1)/8$.
\end{corollary}

% ==================================================
\section{Proof of Theorem~\ref{thm:general}}
\label{sec:proof-general}
% ==================================================

\subsection*{4.1. Active windows and the bound $w(n) \leq \sigma_n$}

\begin{definition}[Active window]
\label{def:active}
An interval $[s, s')$ between consecutive breakpoints of $K^{(n)}$
is \emph{active} if $K^{(n)} \neq 0$ on $(s, s')$.
$w(n)$ denotes the number of active windows.
\end{definition}

\begin{lemma}
\label{lem:w-bound}
$w(n) \leq \sigma_n$.
\end{lemma}

\begin{proof}
By Lemma~\ref{lem:K-vanish} the support of $K^{(n)}$ is contained
in $[0, \sigma_n) \cap \mathbb{Z}$, which contains $\sigma_n$
integers. The number of active windows does not exceed the number
of breakpoints, which is at most $\sigma_n$.
\end{proof}

\begin{lemma}[Parity cancellation]
\label{lem:cancel}
If $\pi_n(s) = 0$, then $s$ is not a breakpoint of $K^{(n)}$,
reducing $w(n)$ below $\sigma_n$.
This occurs when there exist $S_1 \neq S_2$ with
$\sigma(S_1) = \sigma(S_2) = s$ and $|S_1| \not\equiv |S_2| \pmod{2}$.
\end{lemma}

\begin{proof}
By the definition of $\pi_n$: contributions of opposite parity cancel.
\end{proof}

\begin{remark}
\label{rem:cancellation-empirical}
The sparsity ratio
$1 - w(n)/\sigma_n$ grows from $0\%$ at $n=2$
to $100{.}0\%$ at $n \geq 21$ for the sequence of the first primes
(Table~\ref{tab:mertens};
values $100{.}0\%$ result from rounding to one decimal place).
A combinatorial proof of the strict inequality $w(n) < 2^n - 1$
for every $n \geq 3$ remains open
(Section~\ref{sec:open}, item~(iii)).
\end{remark}

\subsection*{4.2. Incremental algorithm}

\begin{lemma}[Window update]
\label{lem:incremental}
The active windows for $n+1$ generators are determined by
the new breakpoints arising from $\pi_{n+1}(s) = \pi_n(s) - \pi_n(s - p_{n+1})$
(Lemma~\ref{lem:parity-recurrence}).
The transition from $n$ to $n+1$ generators costs $O(\sigma_n)$.
\end{lemma}

\begin{proof}
The support of $\pi_n$ is contained in $[0, \sigma_n]$, so
the update $\pi_{n+1}(s) = \pi_n(s) - \pi_n(s - p_{n+1})$
requires at most $\sigma_n + 1$ operations.
\end{proof}

\begin{proof}[Proof of Theorem~\ref{thm:general}]
We compute $\mathcal{W}$ --- the set of triples $(lo, hi, c)$
representing active windows with values of $K^{(n)}$ --- by the
incremental algorithm: start from $n = 1$ (one window $[0, p_1)$
with value $+1$) and perform $n-1$ steps of Lemma~\ref{lem:incremental}.
Total cost: $\sum_{k=1}^{n-1} O(\sigma_k) = O(n \cdot \sigma_n)$.

Then for each window $(lo, hi, c) \in \mathcal{W}$
perform one sliding-window pass over $G$ in time $O(N)$:
accumulate $c \cdot \sum_{j \in G \cap [lo, hi)} n_j$.
Total cost: $O(T_{\mathrm{prep}} + w(n) \cdot N)$.
\end{proof}

\begin{remark}
\label{rem:activation}
The change in $K^{(n)}$ upon adding generator $p_{n+1}$ is
\[
  K^{(n+1)}(d) - K^{(n)}(d) =
  \sum_{\substack{S \subseteq \{p_1,\ldots,p_n\} \\ \sigma(S)+p_{n+1} \leq d}}
  (-1)^{|S|+1}.
\]
A position inactive under $K^{(n)}$ becomes active under $K^{(n+1)}$
if and only if this sum is nonzero.
Empirically $K^{(n)}(d) \notin \{-1, 0, 1\}$ is possible;
values $\pm 3$ are observed at $n = 12$
(Table~\ref{tab:positions}).
\end{remark}

% ==================================================
\section{Proof of Theorem~\ref{thm:density}}
\label{sec:proof-density}
% ==================================================

\begin{proof}[Proof of Theorem~\ref{thm:density}]
We show that $w(n)/((2^n-1) P(n)) \to 0$.

\textbf{Step 1.}
By Lemma~\ref{lem:w-bound}: $w(n) \leq \sigma_n \leq n\,p_n$.

\textbf{Step 2.}
Since $p_1 < p_2 < \cdots < p_n$ are pairwise distinct and $\geq 2$,
we have $p_k \geq k+1$ for each $k$, so
\[
  H_n := \sum_{k=1}^n \frac{1}{p_k} \leq \sum_{j=2}^{n+1} \frac{1}{j} \leq \log(n+1).
\]
From the inequality $\log(1-x) \geq -2x$ for $x \in [0, 1/2]$:
\[
  \log P(n) = \sum_{k=1}^n \log\!\Bigl(1 - \tfrac{1}{p_k}\Bigr)
  \geq -2H_n \geq -2\log(n+1),
\]
whence $P(n) \geq (n+1)^{-2}$.

\textbf{Step 3.}
From Steps 1--2 and $2^n - 1 \geq 2^{n-1}$:
\[
  \frac{w(n)}{(2^n-1)\,P(n)} \leq \frac{2\,n\,p_n\,(n+1)^2}{2^n}.
\]

\textbf{Step 4.}
By the hypothesis $\log p_n = o(n)$:
\[
  \log\frac{p_n}{2^n} = \log p_n - n\log 2 = o(n) - n\log 2 \to -\infty,
\]
so $p_n / 2^n \to 0$, while the factor $2n(n+1)^2$ grows polynomially.
Hence the expression in Step~3 tends to $0$,
which is equivalent to $\delta(n)/P(n) \to 1$.
\end{proof}

\begin{remark}
\label{rem:density-empirical}
For $n \geq 25$ in Table~\ref{tab:mertens}:
$w(n)/(2^n-1) < 10^{-4}$, i.e., $w(n)/(2^n-1)$ is
indistinguishable from zero at four decimal places of $P(n)$
across all eight tested generator families.
\end{remark}

% ==================================================
\section{Open problems}
\label{sec:open}
% ==================================================

\begin{enumerate}[(i)]

\item \textbf{Boundedness of $|K^{(n)}(d)|$.}
Does $\sup_{n,d} |K^{(n)}(d)| < \infty$?
Empirically $|K^{(n)}(d)| \leq 3$ for $n \leq 12$
(Table~\ref{tab:positions}).

\item \textbf{Stability of values.}
For each fixed $d$, does there exist $n_0 = n_0(d)$ such that
$K^{(n)}(d) = K^{(n_0)}(d)$ for all $n \geq n_0$?

\item \textbf{Sharp inequality $w(n) < 2^n - 1$.}
Theorem~\ref{thm:density} proves $w(n)/(2^n-1) \to 0$,
but does not give a sharp bound.
Does $w(n) < 2^n - 1$ hold for every $n \geq 3$
and every sequence of coprime generators?

\item \textbf{Characterisation of $n^*$.}
Let $n^* = \min\{n : \delta(n) > 0\}$.
Theorem~\ref{thm:density} proves $\delta(n)/P(n) \to 1$,
but does not characterise $n^*$.
Conjecture~\ref{conj:harmonic} suggests that $n^*$ depends only
on the harmonic sum $H_n$.

\item \textbf{Monotone behaviour of $\delta(n)$.}
Empirically $\delta(n)$ attains its maximum at $n = 16$
and then decreases (Table~\ref{tab:mertens}).
Does $\delta(n)$ attain a proper maximum at finite $n$?

\item \textbf{Structure of $\Delta w(n)$.}
From $n = 19$ onwards: the increments $\Delta w(n) = w(n) - w(n-1)$
form two increasing subsequences according to the parity of $n$.
What is the arithmetic mechanism behind this split?

\item \textbf{Asymptotics of $\Delta w(n)$.}
Does $\Delta w(n) / p_n \to 1$ for sequences with $p_n \to \infty$?
Numerical data: $p_{300} = 1987$, $\Delta w(300) = 1986$
(baseline sequence); $p_{300} = 2153$, $\Delta w(300) = 2153$
(sequence starting from $101$). The convergence $\Delta w(n) \sim p_n$
would imply $w(n) \sim \sigma_n$.

\end{enumerate}

\begin{conjecture}
\label{conj:harmonic}
Let $n^* = \min\{n : \delta(n) > 0\}$.
The value $n^*$ depends only on the harmonic sum
$H_n = \sum_{k=1}^n 1/p_k$: sequences with the same $H_n$
as the first $n$ primes share the same $n^*$,
regardless of the arithmetic nature of the generators.
\end{conjecture}

The conjecture is confirmed experimentally
(Table~\ref{tab:sequences}).
A formal proof requires the characterisation of $n^*$
described in problem~(iv).

\medskip
\textbf{Related literature.}
Parity cancellation in inclusion--exclusion appears in unrelated contexts:
in the theory of chromatic polynomials
(Whitney~\cite{Whitney1932}; Chen and Qian~\cite{ChenQian2018})
and in analytic sieve theory
(Tao~\cite{Tao2007}).
The interval sparsity of $K^{(n)}$ and parity cancellation
for numerical semigroup kernels have not, to our knowledge,
been described previously.

% ==================================================
\section{Conclusions}
\label{sec:conclusion}
% ==================================================

The starting point is Huang's observation~\cite{Huang2026}:
the quadratic form $Q$ on the gap set of the numerical semigroup
$\langle a, b \rangle$ recovers the combinatorial statistic $\mathrm{dinv}$.
The present paper addresses the computational efficiency of $Q$
and reveals a structure that has not been described before.

\textbf{Complexity reduction.}
The interval structure of $K$ (Lemma~\ref{lem:K-intervals})
allows one to replace the double sum $O(N^2)$ by two sliding-window
passes in time $O(N)$ (Theorem~\ref{thm:linear}).
With $n$ generators the kernel $K^{(n)}$ is defined by
$2^n$ inclusion--exclusion terms, but due to parity cancellation
it has at most $w(n) \leq \sigma_n$ active windows
(Theorem~\ref{thm:general}).
Empirically $w(n) \sim \sigma_n$ while $w(n) \ll 2^n$:
exponential combinatorics reduces to a sparse interval structure,
and the incremental algorithm determines this structure in time
$O(n \cdot \sigma_n)$.

\textbf{Asymptotics.}
Theorem~\ref{thm:density} quantifies this phenomenon:
under the subexponential condition $\log p_n = o(n)$,
\[
  \frac{w(n)}{2^n - 1} = o\!\left(\prod_{i=1}^n
    \Bigl(1 - \frac{1}{p_i}\Bigr)\right), \qquad n \to \infty.
\]
The density of active windows is asymptotically negligible
relative to the density product $P(n)$; equivalently $\delta(n)/P(n) \to 1$.
Numerical data (Table~\ref{tab:mertens}) confirm the convergence
for $n \leq 30$ and indicate a maximum of $\delta(n)$ at $n = 16$,
whose combinatorial characterisation remains open.

\textbf{Outlook.}
Parity cancellation in $K^{(n)}$ admits an algebraic formulation
via the difference operator
$\Delta_{p} f(d) = f(d) - f(d - p)$
acting on the indicator function $\mathbf{1}_{[0,\infty)}$:
\[
  K^{(n)} = \Delta_{p_1} \circ \cdots \circ \Delta_{p_n}\,
             \mathbf{1}_{[0,\infty)}.
\]
In this formulation $\pi_n(s)$ is the coefficient of the operator
mask, and the sparsity of $K^{(n)}$ corresponds to zero coefficients
of that mask. The connection with generating functions, operator
algebra, and harmonic analysis on $\mathbb{Z}$ points to directions
for further research.

% ==================================================
\appendix
\section{Numerical verification}
\label{sec:verification}
% ==================================================

\section*{A.1. Experimental conditions}
\label{subsec:a1-conditions}

Experiments were conducted in Python~3.14 (SymPy~1.14, NumPy)
on a workstation AMD Ryzen~9~5900X (64\,GB, Windows~11~Pro)
and a laptop Intel Core i7-13620H (16\,GB, Windows~11~Pro).
Arithmetic: IEEE~754 double precision ($\varepsilon \approx 2{.}2 \times 10^{-16}$).

Three implementations --- naive $O(N^2)$, prefix sums $O(N\log N)$,
and sliding window $O(N)$ --- form a verification chain.
Operation counts are deterministic functions of $(a, b)$.
Cross-checking by direct enumeration $O(2^n)$
at $n = 27, 28$ was performed multithreaded;
all remaining integer values
($w(n)$, operation counts, active positions)
were bitwise identical on both platforms.
All verification scripts are available at
\url{https://github.com/Flaman55/RelMathApps/tree/main/number_theory/numerical_semigroups/linear_Q}:
\texttt{verify\_linear\_Q.py} (Theorems~\ref{thm:linear}--\ref{thm:general},
Tables~\ref{tab:two-gen}--\ref{tab:multi-gen}),
\texttt{appendix\_verify.py} (Tables~\ref{tab:mertens}--\ref{tab:sequences}),
\texttt{incremental\_windows.py} (incremental algorithm),
and \texttt{section\_C\_parallel.py} (parallel enumeration, $n=19,\ldots,30$).

\section*{A.2. Two generators}
\label{subsec:a2-two}

\begin{table}[H]
\centering
\begin{tabular}{rrrrrr}
\toprule
$a$ & $b$ & $N$ & ops.\ $O(N^2)$ & ops.\ $O(N)$ & max.\ error \\
\midrule
  3 &   5 &     4 &          16 &        16 & $5{.}3\times10^{-15}$ \\
  5 &   7 &    12 &         144 &        48 & $1{.}8\times10^{-14}$ \\
  7 &  11 &    31 &         961 &       124 & $1{.}1\times10^{-13}$ \\
 11 &  13 &    60 &        3600 &       240 & $2{.}8\times10^{-13}$ \\
 23 &  29 &   314 &       98596 &      1256 & $1{.}2\times10^{-11}$ \\
 37 &  41 &   721 &      519841 &      2884 & $4{.}3\times10^{-11}$ \\
 97 & 101 &  4801 &   23049601 &     19204 & $1{.}3\times10^{-9}$  \\
\bottomrule
\multicolumn{6}{l}{\footnotesize All 450 pairs $(a,b)$ with $b \leq 40$: \textbf{PASS}.
  Tolerance $10^{-9}$ ($10^{-8}$ for $N > 500$).}
\end{tabular}
\caption{Verification of Theorem~\ref{thm:linear}: two generators.
Operation ratio $N^2/(4N) = N/4$.}
\label{tab:two-gen}
\end{table}

\section*{A.3. Multiple generators}
\label{subsec:a3-multiple}

\begin{table}[H]
\centering
\begin{tabular}{lrrrr}
\toprule
Generators & $N$ & $w(n)$ & max.\ error & Status \\
\midrule
$\{2,3\}$     & 1 & 2 & $0$                   & OK \\
$\{3,5\}$     & 4 & 2 & $8{.}9\times10^{-16}$ & OK \\
$\{2,3,5\}$   & 1 & 4 & $0$                   & OK \\
$\{3,5,7\}$   & 3 & 5 & $1{.}8\times10^{-15}$ & OK \\
$\{2,3,5,7\}$ & 1 & 8 & $0$                   & OK \\
\bottomrule
\end{tabular}
\caption{Verification of Theorem~\ref{thm:general}: multiple generators
against naive $O(N^2)$.}
\label{tab:multi-gen}
\end{table}

\newpage
\section*{A.4. Density data (Theorem~\ref{thm:density})}
\label{subsec:a4-density}

\begin{table}[H]
\centering
\small
\begin{tabular}{r|rrrrr}
\toprule
$n$ & $2^n-1$ & $w(n)$ & $P(n)$ & Sparsity\% & $\delta(n)$ \\
\midrule
 2 &          3 &    2 & 0{.}3333 &  0{.}0\% & $-0{.}3333$ \\
 3 &          7 &    4 & 0{.}2667 & 14{.}3\% & $-0{.}3048$ \\
 4 &         15 &    8 & 0{.}2286 & 26{.}7\% & $-0{.}3048$ \\
 5 &         31 &   14 & 0{.}2078 & 29{.}0\% & $-0{.}2438$ \\
 6 &         63 &   20 & 0{.}1918 & 44{.}4\% & $-0{.}1257$ \\
 7 &        127 &   30 & 0{.}1805 & 59{.}1\% & $-0{.}0557$ \\
 8 &        255 &   38 & 0{.}1710 & 72{.}2\% & $+0{.}0220$ \\
 9 &        511 &   48 & 0{.}1636 & 81{.}6\% & $+0{.}0697$ \\
10 &       1023 &   60 & 0{.}1579 & 88{.}0\% & $+0{.}0993$ \\
11 &       2047 &   84 & 0{.}1529 & 92{.}5\% & $+0{.}1118$ \\
12 &       4095 &  116 & 0{.}1487 & 95{.}3\% & $+0{.}1204$ \\
13 &       8191 &  146 & 0{.}1451 & 97{.}2\% & $+0{.}1273$ \\
14 &      16383 &  198 & 0{.}1417 & 98{.}3\% & $+0{.}1296$ \\
15 &      32767 &  236 & 0{.}1387 & 99{.}0\% & $+0{.}1315$ \\
16 &      65535 &  288 & 0{.}1361 & 99{.}4\% & $\mathbf{+0{.}1317}$ \\
17 &     131071 &  352 & 0{.}1338 & 99{.}7\% & $+0{.}1311$ \\
18 &     262143 &  412 & 0{.}1316 & 99{.}8\% & $+0{.}1300$ \\
19 &     524287 &  492 & 0{.}1296 & 99{.}9\% & $+0{.}1287$ \\
20 &    1048575 &  564 & 0{.}1278 & 99{.}9\% & $+0{.}1273$ \\
21 &    2097151 &  646 & 0{.}1260 &100{.}0\% & $+0{.}1257$ \\
22 &    4194303 &  720 & 0{.}1245 &100{.}0\% & $+0{.}1243$ \\
23 &    8388607 &  808 & 0{.}1230 &100{.}0\% & $+0{.}1229$ \\
24 &   16777215 &  892 & 0{.}1216 &100{.}0\% & $+0{.}1215$ \\
25 &   33554431 &  996 & 0{.}1203 &100{.}0\% & $+0{.}1203$ \\
26 &   67108863 & 1096 & 0{.}1191 &100{.}0\% & $+0{.}1191$ \\
27 &  134217727 & 1200 & 0{.}1180 &100{.}0\% & $+0{.}1180$ \\
28 &  268435455 & 1306 & 0{.}1169 &100{.}0\% & $+0{.}1169$ \\
29 &  536870911 & 1416 & 0{.}1158 &100{.}0\% & $+0{.}1158$ \\
30 & 1073741823 & 1528 & 0{.}1148 &100{.}0\% & $+0{.}1148$ \\
\bottomrule
\end{tabular}
\caption{Number of active windows $w(n)$, density product $P(n)$,
sparsity $1 - w(n)/\sigma_n$,
and density defect $\delta(n)$ (Theorem~\ref{thm:density})
for $p_i = i$-th prime.
Milestones: $n=8$ --- first $\delta(n) > 0$;
$n=21$ --- sparsity reaches $100\%$ (after rounding);
$n=25$ --- $w(n)/(2^n-1) < 10^{-4}$.
Maximum of $\delta(n)$ at $n=16$ (bold).
Values ``$100.0\%$'' result from rounding to one decimal place.
Data $n \leq 18$: direct enumeration;
$n = 19, \ldots, 30$: incremental algorithm,
independently confirmed for $n=27,28$ by $O(2^n)$ enumeration.}
\label{tab:mertens}
\end{table}

\newpage
\section*{A.5. Values of $K^{(n)}(d)$ --- activation and deactivation}
\label{subsec:a5-values}

\begin{table}[H]
\centering
\small
\begin{tabular}{r|rrrrrrrrrrr}
\toprule
$d$ & $n{=}2$ & $n{=}3$ & $n{=}4$ & $n{=}5$ & $n{=}6$ &
      $n{=}7$ & $n{=}8$ & $n{=}9$ & $n{=}10$ & $n{=}11$ & $n{=}12$ \\
\midrule
 0 &  1 &  1 &  1 &  1 &  1 &  1 &  1 &  1 &  1 &  1 &  1 \\
 3 & -1 & -1 & -1 & -1 & -1 & -1 & -1 & -1 & -1 & -1 & -1 \\
 5 &  $\cdot$ & -1 & -1 & -1 & -1 & -1 & -1 & -1 & -1 & -1 & -1 \\
 7 &  $\cdot$ &  $\cdot$ & -1 & -1 & -1 & -1 & -1 & -1 & -1 & -1 & -1 \\
 8 &  $\cdot$ &  1 &  $\cdot$ &  $\cdot$ &  $\cdot$ &  $\cdot$ &  $\cdot$ &  $\cdot$ &  $\cdot$ &  $\cdot$ &  $\cdot$ \\
 9 &  $\cdot$ &  $\cdot$ &  1 &  1 &  1 &  1 &  1 &  1 &  1 &  1 &  1 \\
10 &  $\cdot$ &  $\cdot$ &  1 &  1 &  1 &  1 &  1 &  1 &  1 &  1 &  1 \\
16 &  $\cdot$ &  $\cdot$ &  $\cdot$ &  $\cdot$ &  1 &  1 &  1 &  1 &  1 &  1 &  1 \\
17 &  $\cdot$ &  $\cdot$ &  $\cdot$ &  1 &  2 &  1 &  1 &  1 &  1 &  1 &  1 \\
25 &  $\cdot$ &  $\cdot$ &  $\cdot$ &  $\cdot$ &  $\cdot$ &  $\cdot$ &  1 &  1 &  1 &  1 &  1 \\
47 &  $\cdot$ &  $\cdot$ &  $\cdot$ &  $\cdot$ &  $\cdot$ &  1 &  1 &  $\cdot$ & -1 & -2 & -3 \\
\bottomrule
\end{tabular}
\caption{Selected values of $K^{(n)}(d)$ for $n = 2, \ldots, 12$
($p_i = i$-th prime). Dot ($\cdot$): $K^{(n)}(d) = 0$.
Position $d = 8$: active at $n=3$, deactivated at $n=4$.
Position $d = 17$: $K^{(6)}(17) = 2$.
Position $d = 47$: $K^{(12)}(47) = -3$.
The table provides empirical examples of values
$|K^{(n)}(d)| = 2, 3$, motivating open problem~(i).
The phenomenon of position activation and deactivation,
empirically arising from parity cancellation
among overlapping subset-sum representations,
is the subject of open problem~(ii).}
\label{tab:positions}
\end{table}


\section*{A.6. Verification of Conjecture~\ref{conj:harmonic}}
\label{subsec:a6-conjecture}

\begin{table}[H]
\centering
\small
\begin{tabular}{lrrrr}
\toprule
Generator sequence & First gen. & $n^*$
  & $\delta/P$, $n=25$ & $\delta/P$, $n=300$ \\
\midrule
\multicolumn{5}{l}{\textit{Prime sequences}} \\
First $n$ primes             &    2 & 8 & 0{.}9998 & 1{.}0000 \\
Primes from $101$            &  101 & 2 & 0{.}9999 & 1{.}0000 \\
Primes from $1009$           & 1009 & 2 & 0{.}9997 & 1{.}0000 \\
Every other prime ($p_2, p_4, \ldots$) &    3 & 8 & 0{.}9998 & 1{.}0000 \\
Every third prime ($p_3, p_6, \ldots$) &  5 & 2 & 0{.}9998 & 1{.}0000 \\
Smaller of twin prime pairs  &    3 & 9 & 0{.}9996 & 1{.}0000 \\
Primes $\equiv 3 \pmod{4}$   &  3 & 8 & 0{.}9998 & 1{.}0000 \\
Primes $\equiv 1 \pmod{4}$   &  5 & 2 & 0{.}9998 & 1{.}0000 \\
\multicolumn{5}{l}{\textit{Random sequences (random odd integers)}} \\
Same $H_n$ as baseline (seed 42)  & {---} & 8 & 0{.}9999 & {---} \\
Same $H_n$ as baseline (seed 137) & {---} & 8 & 0{.}9999 & {---} \\
Different $H_n$ (seed 42)         & {---} & 2 & 0{.}9999 & {---} \\
\bottomrule
\end{tabular}
\caption{Data consistent with Theorem~\ref{thm:density}:
$\delta(n)/P(n) \to 1$ for all tested sequences.
At $n=300$: $\delta(n)/P(n) = 1{.}0000$ (six decimal places)
for all prime sequences.
Random sequences with the same $H_n$ as the baseline yield
identical $n^*$, consistent with Conjecture~\ref{conj:harmonic}.}
\label{tab:sequences}
\end{table}

\clearpage
\printbibliography

\end{document}
